% bismi ALLAH ARRAHMAN ARRAHIM























































\documentclass[11pt,a4paper]{report}

% ============================================================
% EX7 - DOCUMENT DE SYNTHESE DU PROJET
% ============================================================

\usepackage[utf8]{inputenc}
\usepackage[T1]{fontenc}
\usepackage[french]{babel}
\usepackage{lmodern}
\usepackage{microtype}

\usepackage{amsmath,amssymb,mathtools,bm}
\usepackage{booktabs,tabularx,array,multirow}
\usepackage{geometry}
\usepackage{xcolor}
\usepackage{graphicx}
\usepackage{float}
\usepackage{enumitem}
\usepackage{hyperref}
\usepackage{fancyhdr}

\usepackage{tikz}
\usetikzlibrary{
  arrows.meta,
  positioning,
  calc,
  fit,
  backgrounds,
  shapes.geometric,
  shapes.misc
}

\geometry{
  a4paper,
  left=2.15cm,
  right=2.15cm,
  top=2.2cm,
  bottom=2.2cm
}

\hypersetup{
  colorlinks=true,
  linkcolor=blue!55!black,
  urlcolor=blue!55!black,
  citecolor=blue!55!black,
  pdftitle={EX7 - Synthese generale du projet},
  pdfsubject={Architecture modulaire a sondes RBF et matricielles}
}

\pagestyle{fancy}
\fancyhf{}
\fancyhead[L]{Projet EX7}
\fancyhead[R]{Synthèse technique et vision}
\fancyfoot[C]{\thepage}
\setlength{\headheight}{14pt}

\setlength{\parindent}{0pt}
\setlength{\parskip}{5pt}

\setlist[itemize]{leftmargin=1.5em,itemsep=2pt}
\setlist[enumerate]{leftmargin=1.7em,itemsep=2pt}

\definecolor{blueblock}{RGB}{220,235,250}
\definecolor{greenblock}{RGB}{222,242,229}
\definecolor{purpleblock}{RGB}{235,226,248}
\definecolor{orangeblock}{RGB}{250,234,211}
\definecolor{grayblock}{RGB}{240,240,240}
\definecolor{redblock}{RGB}{250,225,225}

\tikzset{
  exblock/.style={
    draw,
    rounded corners=2pt,
    thick,
    align=center,
    text width=3.0cm,
    minimum height=0.85cm,
    inner sep=5pt,
    font=\small
  },
  input/.style={
    exblock,
    fill=grayblock
  },
  local/.style={
    exblock,
    fill=greenblock
  },
  rbf/.style={
    exblock,
    fill=purpleblock
  },
  global/.style={
    exblock,
    fill=orangeblock
  },
  warning/.style={
    exblock,
    fill=redblock
  },
  flow/.style={
    -{Latex[length=2mm]},
    thick
  },
  optionalflow/.style={
    -{Latex[length=2mm]},
    thick,
    dashed
  }
}

\newcommand{\R}{\mathbb{R}}
\newcommand{\softmax}{\operatorname{softmax}}
\newcommand{\ReLU}{\operatorname{ReLU}}
\newcommand{\RBF}{\operatorname{RBF}}
\newcommand{\concat}{\operatorname{concat}}
\newcommand{\vect}[1]{\bm{#1}}

\title{
  \Huge\textbf{EX7}\\[0.3cm]
  \Large Synthèse générale du projet\\[0.2cm]
  \large Architecture modulaire à sondes, coordination RBF,
  calcul matriciel et perspectives linguistiques
}

\author{Document de recherche et de conception}
\date{\today}

\begin{document}

\maketitle

\begin{abstract}
Ce document rassemble la réflexion technique, architecturale
et prospective autour du projet EX7.

Le point de départ est une architecture expérimentale composée
de sept sondes AdaptiveRBF indépendantes et d'un coordinateur
global. À partir de cette expérience, plusieurs évolutions
ont été envisagées : remplacement des calculs RBF locaux par
des transformations matricielles, conservation des RBF dans
un coordinateur global, vectorisation des sondes, distribution
du calcul entre unités spécialisées et extension possible
vers le traitement du langage naturel.

La vision à long terme est d'explorer une architecture dans
laquelle des unités de traitement relativement petites
coopèrent pour produire des représentations complexes,
modéliser des relations et traiter des séquences linguistiques.

Cette vision est une hypothèse de recherche, non un résultat
déjà démontré. Le document distingue donc systématiquement
l'implémentation existante, les propositions architecturales
et les questions scientifiques ouvertes.
\end{abstract}

\tableofcontents
\listoffigures
\listoftables
\newpage

% ============================================================
\chapter{Origine et philosophie du projet}
% ============================================================

\section{Point de départ}

Le projet EX7 part d'une question architecturale :

\begin{quote}
Est-il possible de construire un système d'apprentissage
composé de plusieurs unités spécialisées, capables de traiter
des informations localement puis de coopérer au moyen d'un
mécanisme de coordination global ?
\end{quote}

L'objectif n'est pas simplement de remplacer un réseau
de neurones classique par plusieurs réseaux plus petits.

L'idée centrale est de déplacer une partie de la complexité
du modèle vers son organisation :

\begin{itemize}
  \item plusieurs unités locales ;
  \item des paramètres et des transformations propres
  à chaque unité ;
  \item une éventuelle spécialisation fonctionnelle ;
  \item une coordination des informations produites ;
  \item une possibilité d'évolution indépendante des unités.
\end{itemize}

La nouveauté potentielle peut donc résider dans l'organisation
et la coordination de composants connus, et non nécessairement
dans l'invention d'une nouvelle fonction mathématique élémentaire.

\section{Une distinction fondamentale}

Il faut distinguer trois niveaux :

\begin{enumerate}
  \item \textbf{Le composant :} une fonction de calcul locale,
  par exemple une unité RBF ou un petit bloc matriciel.

  \item \textbf{L'organisation :} la manière dont plusieurs
  composants sont distribués, entraînés et connectés.

  \item \textbf{Le comportement global :} les capacités
  émergentes ou apprises par leur coopération.
\end{enumerate}

Une fonction connue peut participer à une architecture
nouvelle. Toutefois, l'intérêt scientifique de cette
organisation doit être démontré par des comparaisons
expérimentales et par rapport aux architectures existantes.

\section{La notion de sonde}

Dans le projet, une sonde désigne une unité de traitement
qui reçoit des informations, effectue une transformation
et produit une représentation ou une contribution.

Une sonde n'est donc pas nécessairement un neurone scalaire
classique.

Elle peut être un module relativement autonome, composé
de plusieurs opérations et de plusieurs paramètres.

Dans EX7 V003, une sonde est un petit expert vectoriel
à base de fonctions radiales, et non une unité élémentaire
équivalente à un seul neurone d'un MLP.

\begin{figure}[H]
\centering
\begin{tikzpicture}[node distance=0.65cm]
  \node[input] (x) {Entrée\\Informations};
  \node[local,below=of x,text width=7cm] (p) {
    Sonde\\
    Transformation locale\\
    Extraction d'une représentation
  };
  \node[global,below=of p] (g) {
    Coordination\\
    Combinaison des résultats
  };
  \node[input,below=of g] (y) {Sortie globale};

  \draw[flow] (x) -- (p);
  \draw[flow] (p) -- (g);
  \draw[flow] (g) -- (y);
\end{tikzpicture}
\caption{Principe général d'une unité de traitement et de sa coordination.}
\end{figure}

% ============================================================
\chapter{EX7 V003 : architecture effectivement implémentée}
% ============================================================

\section{Vue d'ensemble}

EX7 V003 est désigné \textit{OuterCoord-Lite-7x1}.

Il comprend :

\begin{itemize}
  \item sept sondes AdaptiveRBF indépendantes ;
  \item un coordinateur global ;
  \item un mécanisme de combinaison des sorties ;
  \item un prétraitement des caractéristiques.
\end{itemize}

Le problème de référence utilise le jeu Wisconsin Diagnostic
Breast Cancer :

\[
D=30,\qquad C=2,\qquad M=7.
\]

Le nombre de caractéristiques est \(D=30\), le nombre
de classes est \(C=2\), et le nombre de sondes est \(M=7\).

\section{Architecture générale}

\begin{figure}[H]
\centering
\begin{tikzpicture}[node distance=0.65cm]
  \node[input] (x) {
    Entrée standardisée\\
    \(\vect{x}\in\R^{30}\)
  };

  \node[local,below=of x,text width=9.4cm] (p) {
    Sept sondes AdaptiveRBF indépendantes\\[2pt]
    \(\RBF_1(\vect{x}),\ldots,\RBF_7(\vect{x})\)
  };

  \node[global,below=of p,text width=6cm] (rho) {
    Coordinateur global\\
    \(\vect{\rho}=\softmax(\vect{a})\)
  };

  \node[input,below=of rho] (out) {
    Logits de classification\\
    \(\vect{y}\in\R^2\)
  };

  \draw[flow] (x) -- (p);
  \draw[flow] (p) -- (rho);
  \draw[flow] (rho) -- (out);
\end{tikzpicture}
\caption{Architecture conceptuelle de EX7 V003.}
\end{figure}

\section{Structure interne d'une sonde}

Chaque sonde reçoit les 30 caractéristiques d'entrée.

Elle possède un encodeur appris :

\[
\vect{z}_i
=
\tanh(W_i\vect{x}+\vect{b}_i).
\]

Avec :

\[
W_i\in\R^{4\times 30},
\qquad
\vect{z}_i\in\R^4.
\]

Chaque sonde utilise huit centres RBF :

\[
\vect{c}_{i,k}\in\R^4,
\qquad k=1,\ldots,8.
\]

La distance au centre \(k\) est :

\[
d_{i,k}^{2}
=
\|\vect{z}_i-\vect{c}_{i,k}\|_2^2.
\]

La largeur radiale est :

\[
s_{i,k}=\exp(\ell_{i,k}).
\]

L'activation radiale est :

\[
\phi_{i,k}(\vect{x})
=
\exp\left(
-\frac{d_{i,k}^{2}}{2s_{i,k}^{2}}
\right).
\]

La sonde produit ensuite ses logits :

\[
\vect{q}_i(\vect{x})
=
A_i^\top\vect{\phi}_i(\vect{x})
+\vect{c}_i^{out}.
\]

\section{Coordination globale}

Le coordinateur utilise un vecteur de poids appris :

\[
\vect{\rho}
=
\softmax(\vect{a}).
\]

La sortie globale est :

\[
\boxed{
\vect{y}(\vect{x})
=
\sum_{i=1}^{7}
\rho_i\vect{q}_i(\vect{x})
}
\]

Une caractéristique importante est que \(\vect{\rho}\)
est indépendant de l'entrée.

Le coordinateur ne choisit donc pas dynamiquement une sonde
différente selon chaque exemple. Il applique une pondération
globale apprise.

\section{Organisation des données d'entraînement}

Les exemples d'entraînement sont répartis entre les sept
sondes au moyen de partitions stratifiées disjointes.

La répartition concerne les exemples, et non les dimensions
d'entrée.

Chaque sonde reçoit néanmoins les 30 caractéristiques
complètes des exemples qui lui sont attribués pendant
son entraînement.

À l'inférence, toutes les sondes sont exécutées sur chaque
nouvel exemple.

Il faut donc distinguer :

\begin{itemize}
  \item spécialisation par partition des exemples
  pendant l'entraînement ;
  \item exécution de toutes les sondes pendant l'inférence ;
  \item absence de partition dimensionnelle dans V003.
\end{itemize}

\section{Paramètres}

Les chiffres de référence sont :

\[
P_{\mathrm{train}}=1281,
\]

\[
P_{\mathrm{préprocesseur}}=60,
\]

\[
P_{\mathrm{total}}=1341.
\]

Les comparateurs historiques mentionnés sont :

\begin{itemize}
  \item MLP apparié : 1349 paramètres ;
  \item MLP large : 5375 paramètres ;
  \item Orchestra historique : 6712 paramètres.
\end{itemize}

Le comptage des paramètres ne doit pas être confondu
avec le coût total d'exécution ou la consommation mémoire.

% ============================================================
\chapter{Résultats expérimentaux et interprétation}
% ============================================================

\section{Résultats disponibles}

Les résultats communiqués pour les expériences de référence
sont résumés ci-dessous.

\begin{table}[H]
\centering
\small
\begin{tabularx}{\textwidth}{
  >{\raggedright\arraybackslash}X
  r r r
}
\toprule
Modèle & Accuracy & Log-loss & Paramètres\\
\midrule
EX7 V003 & 97,88 \% & 0,06806 & 1341\\
Orchestra OuterCoord-only & 98,59 \% & 0,06013 & 6712\\
Orchestra récursive & 98,35 \% & 0,06187 & 6712\\
MLP apparié & 98,12 \% & 0,05136 & 1349\\
MLP large & 98,59 \% & 0,05222 & 5375\\
\bottomrule
\end{tabularx}
\caption{Résultats de référence communiqués pour les expériences EX7.}
\end{table}

Ces résultats montrent qu'EX7 atteint une accuracy proche
de celle du MLP apparié, avec un nombre de paramètres
comparable.

Cependant, le MLP apparié possède une log-loss plus faible.
L'accuracy seule ne suffit donc pas à conclure à une
équivalence de qualité prédictive.

\section{Latence}

Les latences p50 communiquées sont :

\begin{table}[H]
\centering
\small
\begin{tabularx}{\textwidth}{X r}
\toprule
Modèle & Latence p50\\
\midrule
EX7 V003 & \(0,290\ \mathrm{ms}\)\\
Orchestra & \(1,49\ \mathrm{ms}\) environ\\
MLP apparié & \(0,025\ \mathrm{ms}\)\\
MLP large & \(0,026\ \mathrm{ms}\)\\
\bottomrule
\end{tabularx}
\caption{Latences p50 communiquées.}
\end{table}

EX7 est donc environ :

\[
\frac{0,290}{0,025}=11,6
\]

fois plus lent que le MLP apparié sur cette mesure.

Il est cependant nettement plus rapide que l'Orchestra
historique dans le benchmark considéré.

\section{Pourquoi EX7 est-il plus lent ?}

Le coût d'EX7 vient notamment de la succession de nombreuses
petites opérations :

\begin{enumerate}
  \item sept encodeurs indépendants ;
  \item sept calculs de distances ;
  \item sept ensembles d'activations exponentielles ;
  \item sept lectures RBF ;
  \item coordination finale.
\end{enumerate}

Un MLP réalise principalement des opérations matricielles
et des additions, particulièrement optimisées dans les
bibliothèques de calcul.

Le nombre de paramètres similaire ne garantit donc pas
une latence similaire.

\section{Interprétation équilibrée}

Les résultats permettent de retenir plusieurs constats :

\begin{itemize}
  \item EX7 est une architecture compacte.
  \item Sa performance est proche de celle du MLP apparié
  en accuracy.
  \item Le MLP apparié est meilleur en log-loss dans les
  résultats communiqués.
  \item EX7 est nettement plus lent que le MLP apparié.
  \item EX7 est plus rapide que l'Orchestra historique
  dans le benchmark communiqué.
\end{itemize}

Ces observations justifient des optimisations et de nouvelles
expériences, mais ne démontrent pas une supériorité générale
de l'architecture à sondes.

\section{Limites des résultats}

Les résultats disponibles ne démontrent pas :

\begin{itemize}
  \item une généralisation à des jeux de données externes ;
  \item une supériorité statistique robuste ;
  \item une interprétabilité causale ;
  \item une efficacité énergétique ;
  \item une meilleure efficacité mémoire totale ;
  \item une aptitude clinique au déploiement.
\end{itemize}

L'audit indépendant de 91 contrôles mentionné dans le projet
renforce la cohérence des artefacts examinés, mais ne remplace
pas une validation externe indépendante.

Le problème médical utilisé est ici un banc d'essai
expérimental. Aucun usage clinique ne peut être déduit
de ces résultats.

% ============================================================
\chapter{Le problème de la latence et EX7-Matrix}
% ============================================================

\section{Constat}

La question soulevée est la suivante :

\begin{quote}
Pourquoi ne pas modifier EX7 pour que les sondes utilisent
principalement des multiplications matricielles et des
additions, comme un MLP ?
\end{quote}

Cette proposition vise à conserver l'organisation modulaire
tout en simplifiant les opérations locales.

\section{Sonde matricielle}

Une sonde matricielle peut être définie par :

\[
\vect{h}_i
=
\sigma(W_i\vect{x}+\vect{b}_i),
\]

avec :

\[
W_i\in\R^{d_h\times D}.
\]

Une projection de sortie peut être ajoutée :

\[
\vect{q}_i
=
V_i\vect{h}_i+\vect{c}_i.
\]

Une fonction d'activation non linéaire reste nécessaire.

En effet, une composition d'opérations affines reste affine :

\[
W_2(W_1\vect{x}+\vect{b}_1)+\vect{b}_2
=
W'\vect{x}+\vect{b}'.
\]

ReLU est un choix possible :

\[
\ReLU(u)=\max(0,u).
\]

\section{Architecture EX7-Matrix}

\begin{figure}[H]
\centering
\begin{tikzpicture}[node distance=0.65cm]
  \node[input] (x) {Entrée complète};
  \node[local,below left=1cm and 0.2cm of x] (p1) {
    Sonde 1\\
    Matrices + ReLU
  };
  \node[local,below=1cm of x] (p2) {
    Sonde 2\\
    Matrices + ReLU
  };
  \node[local,below right=1cm and 0.2cm of x] (p7) {
    Sonde 7\\
    Matrices + ReLU
  };
  \node[global,below=1.1cm of p2,text width=5cm] (g) {
    Coordinateur\\
    Fusion des contributions
  };
  \node[input,below=of g] (y) {Sortie};

  \draw[flow] (x) -- (p1);
  \draw[flow] (x) -- (p2);
  \draw[flow] (x) -- (p7);
  \draw[flow] (p1) -- (g);
  \draw[flow] (p2) -- (g);
  \draw[flow] (p7) -- (g);
  \draw[flow] (g) -- (y);
\end{tikzpicture}
\caption{Principe d'EX7-Matrix.}
\end{figure}

Les sondes restent indépendantes, mais leur moteur local
est matriciel.

\section{Coordination matricielle}

Une première variante conserve le principe de pondération
globale :

\[
\vect{\rho}=\softmax(\vect{a}).
\]

La sortie devient :

\[
\boxed{
\vect{y}
=
\sum_{i=1}^{M}
\rho_i
\left[
V_i\sigma(W_i\vect{x}+\vect{b}_i)
+\vect{c}_i
\right]
}
\]

Cette variante permet d'isoler l'effet du remplacement
des calculs RBF locaux.

\section{Vectorisation}

Les sept sondes peuvent être regroupées dans un tenseur
de poids :

\[
\mathcal{W}
\in\R^{M\times d_h\times D}.
\]

Pour un batch :

\[
X\in\R^{B\times D},
\]

le calcul peut être organisé sous forme d'opérations
matricielles batchées.

La vectorisation peut réduire le coût d'orchestration
des petites opérations, mais son efficacité dépend
du matériel et de la taille des batchs.

Il faut mesurer séparément :

\begin{itemize}
  \item batch de taille 1 ;
  \item batchs de tailles intermédiaires ;
  \item batchs plus grands ;
  \item latences p50, p95 et p99.
\end{itemize}

La vectorisation est une optimisation d'implémentation,
pas nécessairement une nouvelle architecture mathématique.

% ============================================================
\chapter{Architecture hybride : RBF global et Matrix local}
% ============================================================

\section{Origine de la proposition}

Une question importante a émergé :

\begin{quote}
Pourquoi choisir entre EX7 original et une architecture
matricielle qui abandonnerait les RBF ?
Ne pourrait-on pas transférer la fonction RBF vers la
coordination globale et conserver des sondes matricielles
locales ?
\end{quote}

Cette proposition constitue une troisième voie.

Elle conserve :
\begin{itemize}
  \item les sondes indépendantes ;
  \item les transformations locales rapides ;
  \item la non-linéarité RBF ;
  \item une coordination globale.
\end{itemize}

La différence est que les RBF ne sont plus calculées
dans chaque sonde, mais dans un module global.

\section{Définition de EX7-Hybrid}

Chaque sonde calcule :

\[
\vect{h}_i
=
\sigma(W_i\vect{x}+\vect{b}_i).
\]

Les représentations sont concaténées :

\[
\vect{H}
=
\concat(\vect{h}_1,\ldots,\vect{h}_M).
\]

On peut ensuite utiliser une projection globale :

\[
\vect{u}=P\vect{H}+\vect{b}_P.
\]

Le coordinateur RBF reçoit \(\vect{u}\).

Pour chaque centre global \(\vect{\mu}_k\) :

\[
\psi_k(\vect{u})
=
\exp\left(
-\frac{
\|\vect{u}-\vect{\mu}_k\|_2^2
}{
2\tau_k^2
}
\right).
\]

La sortie finale est :

\[
\boxed{
\vect{y}
=
A_g^\top\vect{\psi}(\vect{u})
+\vect{b}_g
}
\]

\section{Schéma de l'architecture hybride}

\begin{figure}[H]
\centering
\begin{tikzpicture}[node distance=0.62cm]
  \node[input] (x) {Entrée};
  \node[local,below=of x,text width=8cm] (local) {
    Sept sondes matricielles\\
    \(\vect{h}_i=\sigma(W_i\vect{x}+\vect{b}_i)\)
  };
  \node[local,below=of local,text width=8cm] (fusion) {
    Fusion des représentations\\
    \(\vect{H}=\concat(\vect{h}_1,\ldots,\vect{h}_7)\)
  };
  \node[rbf,below=of fusion,text width=8cm] (rbf) {
    Coordinateur RBF global\\
    Centres, distances, largeurs et activations
  };
  \node[input,below=of rbf] (out) {Prédiction};

  \draw[flow] (x) -- (local);
  \draw[flow] (local) -- (fusion);
  \draw[flow] (fusion) -- (rbf);
  \draw[flow] (rbf) -- (out);
\end{tikzpicture}
\caption{Architecture hybride proposée.}
\end{figure}

\section{Hypothèse scientifique}

L'hypothèse est que les RBF pourraient être particulièrement
utiles pour reconnaître des configurations globales dans
les représentations produites localement.

Cela permettrait de séparer :

\begin{itemize}
  \item l'extraction locale de caractéristiques ;
  \item la reconnaissance de configurations globales ;
  \item la production de la sortie.
\end{itemize}

Ce bénéfice n'est pas démontré. Le coordinateur global
pourrait devenir coûteux si sa dimension ou son nombre
de centres est trop important.

\section{Comparaison des quatre variantes}

\begin{table}[H]
\centering
\small
\begin{tabularx}{\textwidth}{
  >{\raggedright\arraybackslash}p{2.4cm}
  X
  >{\raggedright\arraybackslash}p{2.4cm}
}
\toprule
Variante & Calcul local & Coordination\\
\midrule
EX7-RBF
&
Encodeur et fonctions RBF
&
Pondération globale\\

EX7-Matrix
&
Matrices et activation
&
Pondération globale\\

EX7-Matrix-Vector
&
Matrices batchées
&
Pondération globale\\

EX7-Hybrid
&
Matrices et activation
&
RBF globale\\
\bottomrule
\end{tabularx}
\caption{Les quatre architectures à comparer.}
\end{table}

% ============================================================
\chapter{Répartition du travail entre les sondes}
% ============================================================

\section{La question des dimensions}

Une autre idée discutée est la possibilité de ne pas demander
à chaque sonde de traiter l'intégralité d'une grande
représentation.

Pour un vecteur de dimension :

\[
D=4096,
\]

on pourrait imaginer :

\[
M=16,\qquad d_{\mathrm{local}}=256.
\]

Chaque sonde recevrait alors un sous-espace différent.

\section{Distinction avec EX7 V003}

Cette organisation n'existe pas dans EX7 V003.

Dans V003 :

\[
\vect{x}_i=\vect{x}
\qquad
\forall i.
\]

Chaque sonde reçoit les 30 caractéristiques complètes.

La spécialisation actuelle concerne les partitions d'exemples
d'entraînement, pas les dimensions.

La distribution dimensionnelle est donc une proposition
future, distincte de l'implémentation de référence.

\section{Proposition de partition dimensionnelle}

On peut écrire :

\[
\vect{x}
=
\begin{bmatrix}
\vect{x}_1\\
\vect{x}_2\\
\vdots\\
\vect{x}_M
\end{bmatrix}.
\]

Chaque sonde traite son sous-espace :

\[
\vect{h}_i
=
f_i(\vect{x}_i).
\]

Le système fusionne ensuite les représentations :

\[
\vect{H}
=
\concat(\vect{h}_1,\ldots,\vect{h}_M).
\]

\begin{figure}[H]
\centering
\begin{tikzpicture}[node distance=0.65cm]
  \node[input] (x) {
    Vecteur global\\
    \(4096\) dimensions
  };

  \node[local,below left=1cm and 0.3cm of x] (p1) {
    Sonde 1\\
    256 dimensions
  };
  \node[local,below=1cm of x] (p2) {
    Sonde 2\\
    256 dimensions
  };
  \node[local,below right=1cm and 0.3cm of x] (p3) {
    Sonde 16\\
    256 dimensions
  };

  \node[global,below=1.15cm of p2,text width=5.4cm] (g) {
    Fusion des représentations
  };

  \draw[flow] (x) -- (p1);
  \draw[flow] (x) -- (p2);
  \draw[flow] (x) -- (p3);
  \draw[flow] (p1) -- (g);
  \draw[flow] (p2) -- (g);
  \draw[flow] (p3) -- (g);
\end{tikzpicture}
\caption{Exemple conceptuel de distribution des dimensions.}
\end{figure}

Le schéma ne représente que trois sondes sur seize
pour conserver sa lisibilité.

\section{Le problème des interactions}

Un découpage strict peut faire perdre des relations
entre dimensions différentes.

Par exemple, si deux informations importantes sont situées
dans deux sous-espaces distincts, une sonde isolée ne peut
pas les combiner directement.

Plusieurs solutions sont envisageables :

\begin{enumerate}
  \item sous-espaces chevauchants ;
  \item projections apprises ;
  \item échanges de messages entre sondes ;
  \item coordinateur global ;
  \item coordination hiérarchique ;
  \item traitement itératif avec retours d'information.
\end{enumerate}

La distribution du calcul ne doit donc pas devenir une
fragmentation qui détruit les relations utiles.

\section{Principe général}

Le véritable objectif n'est pas simplement de réduire
la dimension d'entrée de chaque sonde.

Il est de répartir le travail tout en préservant
la possibilité de recombiner les informations.

La complexité peut être distribuée entre les unités,
mais les interactions entre leurs résultats doivent
rester accessibles au système.

% ============================================================
\chapter{Coordination, communication et spécialisation}
% ============================================================

\section{Du coordinateur fixe à la coordination dynamique}

Le coordinateur EX7 V003 utilise :

\[
\vect{\rho}=\softmax(\vect{a}).
\]

Ce mécanisme attribue des poids globaux aux sondes.

Il ne réalise pas de communication conditionnelle entre
les sondes.

Une évolution possible serait de rendre les poids
dépendants de l'entrée :

\[
\vect{\rho}(\vect{x})
=
\softmax(g(\vect{H})).
\]

La sortie deviendrait :

\[
\vect{y}
=
\sum_{i=1}^{M}
\rho_i(\vect{x})\vect{q}_i.
\]

Cette variante introduirait un mécanisme de routage
ou de pondération dynamique.

Elle doit être distinguée du coordinateur fixe actuel.

\section{Communication entre sondes}

On peut également envisager que chaque sonde reçoive
des informations produites par les autres.

Une formulation générique est :

\[
\vect{m}_i^{(t)}
=
\sum_{j\neq i}
\alpha_{ij}^{(t)}
\,g_{ij}(\vect{h}_j^{(t)}).
\]

La sonde met alors à jour son état :

\[
\vect{h}_i^{(t+1)}
=
f_i\left(
\vect{h}_i^{(t)},
\vect{m}_i^{(t)}
\right).
\]

Cette organisation est un mécanisme de passage de messages.

Elle permettrait à des unités locales de combiner
progressivement des informations provenant d'autres unités.

\section{Spécialisation fonctionnelle}

Il est tentant d'attribuer des rôles explicites aux sondes :

\begin{itemize}
  \item relations lexicales ;
  \item relations syntaxiques ;
  \item relations sémantiques ;
  \item relations contextuelles.
\end{itemize}

Mais ces rôles sont des hypothèses de conception.

Une architecture ne devient pas spécialisée simplement
parce que ses composants portent des noms différents.

Il faut vérifier que les sondes apprennent effectivement
des fonctions distinctes.

On peut utiliser :

\begin{itemize}
  \item des ablations individuelles ;
  \item des permutations de sondes ;
  \item des perturbations ciblées ;
  \item des mesures de similarité des représentations ;
  \item des tests sur des tâches spécialisées.
\end{itemize}

% ============================================================
\chapter{Vers le langage naturel et une architecture linguistique}
% ============================================================

\section{La vision à long terme}

La discussion a progressivement conduit à une ambition
plus large :

\begin{quote}
Construire un modèle capable de traiter le langage naturel
par coopération d'un grand nombre de sondes relativement
petites, plutôt que de concentrer toute la capacité de
traitement dans un ensemble de blocs denses traditionnels.
\end{quote}

Dans cette vision, les sondes ne seraient plus uniquement
des experts de classification.

Elles pourraient participer à la construction de
représentations linguistiques, à l'analyse de relations
et à la prédiction du token suivant.

Cette hypothèse ne signifie pas que le système comprendrait
automatiquement le langage par le seul fait de posséder
beaucoup de sondes.

Il faut une architecture et un objectif d'apprentissage
capables de traiter les séquences, les dépendances
et le contexte.

\section{Le passage de la classification au langage}

EX7 V003 traite un vecteur de caractéristiques :

\[
\vect{x}\in\R^{30}.
\]

Un modèle linguistique doit traiter une séquence :

\[
X=(x_1,x_2,\ldots,x_T),
\]

où chaque \(x_t\) représente un token ou une représentation
de token.

Le système doit prendre en compte :

\begin{itemize}
  \item l'ordre des tokens ;
  \item leur position ;
  \item les dépendances locales ;
  \item les dépendances éloignées ;
  \item les ambiguïtés ;
  \item les relations entre concepts ;
  \item le contexte accumulé.
\end{itemize}

Il faut donc ajouter une architecture séquentielle,
un mécanisme de mémoire ou une autre méthode permettant
de conserver et de transmettre le contexte.

\section{Objectif d'apprentissage possible}

Pour un modèle autorégressif, l'objectif classique
est de prédire le token suivant :

\[
P(x_{t+1}\mid x_1,\ldots,x_t).
\]

La perte moyenne de prédiction peut être écrite :

\[
\mathcal{L}_{\mathrm{LM}}
=
-\frac{1}{T-1}
\sum_{t=1}^{T-1}
\log P(x_{t+1}\mid x_1,\ldots,x_t).
\]

Une architecture à sondes pourrait être entraînée avec
cet objectif.

La difficulté ne réside pas uniquement dans la sortie
du token suivant : elle réside dans la construction
de représentations qui conservent les relations utiles
du contexte.

\section{Architecture linguistique conceptuelle}

\begin{figure}[H]
\centering
\begin{tikzpicture}[node distance=0.65cm]
  \node[input] (tokens) {
    Séquence de tokens\\
    \(x_1,\ldots,x_T\)
  };

  \node[local,below=of tokens,text width=8.5cm] (embed) {
    Représentations d'entrée\\
    Embeddings et informations de position
  };

  \node[local,below=of embed,text width=8.5cm] (probes) {
    Groupe de sondes\\
    Transformations locales et représentations
  };

  \node[global,below=of probes,text width=8.5cm] (coord) {
    Coordination sémantique\\
    Échanges et combinaison des informations
  };

  \node[rbf,below=of coord,text width=8.5cm] (memory) {
    Mémoire et contexte\\
    Représentation actualisée de la séquence
  };

  \node[input,below=of memory] (pred) {
    Prédiction du token suivant
  };

  \draw[flow] (tokens) -- (embed);
  \draw[flow] (embed) -- (probes);
  \draw[flow] (probes) -- (coord);
  \draw[flow] (coord) -- (memory);
  \draw[flow] (memory) -- (pred);
\end{tikzpicture}
\caption{Vision conceptuelle d'un modèle linguistique à sondes.}
\end{figure}

Ce schéma représente une piste de conception, et non
un modèle déjà implémenté.

% ============================================================
\chapter{L'attention est-elle nécessaire ?}
% ============================================================

\section{La question}

Un modèle linguistique doit établir des relations entre
les éléments d'une séquence.

La question n'est pas obligatoirement de conserver
l'attention standard, mais de savoir comment le modèle
va représenter et exploiter ces relations.

\section{Fonction de l'attention}

Dans un Transformer, l'attention utilise des projections
Query, Key et Value :

\[
Q=XW_Q,\qquad K=XW_K,\qquad V=XW_V.
\]

Le calcul standard est :

\[
\operatorname{Attention}(Q,K,V)
=
\softmax\left(
\frac{QK^\top}{\sqrt{d_k}}
\right)V.
\]

Ce mécanisme permet à chaque position de combiner
des informations provenant d'autres positions.

Il est efficace, mais il n'est pas le seul mécanisme
possible de communication contextuelle.

\section{Trois stratégies envisagées}

\subsection{Option A : conserver l'attention}

On conserve l'attention du Transformer et on remplace
les blocs MLP par des sondes.

\begin{figure}[H]
\centering
\begin{tikzpicture}[node distance=0.55cm]
  \node[input] (x) {Tokens};
  \node[global,below=of x,text width=6cm] (att) {
    Attention standard
  };
  \node[local,below=of att,text width=6cm] (p) {
    Sondes matricielles
  };
  \node[input,below=of p] (y) {Nouvelle représentation};

  \draw[flow] (x) -- (att);
  \draw[flow] (att) -- (p);
  \draw[flow] (p) -- (y);
\end{tikzpicture}
\caption{Transformer hybride avec sondes.}
\end{figure}

Cette approche est la plus simple pour mesurer
l'apport des sondes sans modifier simultanément
le mécanisme relationnel.

\subsection{Option B : remplacer l'attention}

Les sondes pourraient communiquer entre elles
au moyen d'un mécanisme de passage de messages.

\[
\vect{m}_i
=
\sum_j
\alpha_{ij}g(\vect{h}_j).
\]

Les mises à jour seraient :

\[
\vect{h}_i'
=
f_i(\vect{h}_i,\vect{m}_i).
\]

Cette voie pourrait remplacer l'attention si les échanges
permettent de capturer les dépendances pertinentes.

\subsection{Option C : mémoire et coordination récurrente}

Une autre possibilité consiste à maintenir un état global :

\[
\vect{s}_{t+1}
=
F(\vect{s}_t,\vect{h}_t).
\]

La mémoire évolue au fil de la séquence et fournit
un contexte aux sondes suivantes.

Cette solution peut éviter de calculer une attention
complète entre toutes les positions, mais elle doit
éviter de perdre les informations importantes du passé.

\section{Conclusion sur l'attention}

L'attention standard n'est pas une nécessité théorique
absolue.

En revanche, le modèle doit fournir une fonction équivalente
ou complémentaire permettant :

\begin{itemize}
  \item de relier des éléments éloignés ;
  \item de gérer les ambiguïtés contextuelles ;
  \item de transmettre des informations entre positions ;
  \item de préserver les relations importantes ;
  \item de construire une représentation cohérente.
\end{itemize}

Supprimer l'attention sans remplacer ses fonctions
ne suffit pas.

% ============================================================
\chapter{Interprétabilité et compréhension des mécanismes}
% ============================================================

\section{Hypothèse}

Une architecture modulaire pourrait faciliter l'analyse
du fonctionnement interne d'un modèle.

Au lieu d'observer un très grand bloc opaque, on pourrait
examiner :

\begin{itemize}
  \item les entrées de chaque sonde ;
  \item ses transformations ;
  \item ses activations ;
  \item ses sorties ;
  \item sa contribution à la prédiction ;
  \item ses interactions avec les autres sondes.
\end{itemize}

Cela pourrait rendre l'analyse plus structurée.

\section{Modularité et interprétabilité}

Il faut néanmoins distinguer modularité et interprétabilité.

Un système composé de modules distincts n'est pas
automatiquement interprétable.

Une sonde peut elle-même apprendre une représentation
complexe et difficile à expliquer.

Il faut donc mesurer si les modules ont réellement
des fonctions identifiables et stables.

\section{Méthodes de vérification}

\begin{enumerate}
  \item \textbf{Ablation :} retirer une sonde et mesurer
  la dégradation du résultat.

  \item \textbf{Intervention :} modifier une représentation
  locale et observer les conséquences.

  \item \textbf{Permutation :} échanger les sorties de deux
  sondes pour vérifier si leurs rôles sont différents.

  \item \textbf{Analyse des activations :} étudier les
  représentations produites pour des catégories de données.

  \item \textbf{Tests contrefactuels :} modifier une partie
  de l'entrée en maintenant le reste constant.

  \item \textbf{Analyse de communication :} mesurer quelles
  informations circulent entre les unités.
\end{enumerate}

\section{Vers une boîte noire partiellement inspectable}

La vision est de construire un système dont le fonctionnement
peut être étudié à plusieurs niveaux :

\begin{itemize}
  \item niveau local : fonctionnement d'une sonde ;
  \item niveau intermédiaire : coopération entre sondes ;
  \item niveau global : coordination et décision ;
  \item niveau séquentiel : évolution du contexte.
\end{itemize}

Cela pourrait améliorer la traçabilité, mais ne garantit
pas à lui seul une explication causale ou complète.

% ============================================================
\chapter{Limites conceptuelles et risques}
% ============================================================

\section{Une sonde ne remplace pas automatiquement un neurone}

Une sonde EX7 est un module comportant plusieurs paramètres
et plusieurs opérations.

Dire qu'une sonde remplace un neurone peut être trompeur.

Il faut distinguer :

\begin{itemize}
  \item remplacement d'un neurone scalaire ;
  \item remplacement d'un petit bloc MLP ;
  \item remplacement d'un bloc fonctionnel d'un Transformer ;
  \item remplacement d'un mécanisme de communication.
\end{itemize}

La portée de l'équivalence doit être définie
par l'expérience et par la fonction calculée.

\section{Un grand nombre de sondes ne garantit pas l'intelligence}

Multiplier les unités ne garantit pas :

\begin{itemize}
  \item la spécialisation ;
  \item la généralisation ;
  \item la compréhension sémantique ;
  \item le raisonnement ;
  \item l'efficacité ;
  \item la stabilité.
\end{itemize}

Ces propriétés doivent être apprises et évaluées.

\section{Le coût de communication}

Répartir le calcul entre de petites unités peut réduire
la taille des opérations locales.

Mais il peut aussi introduire :

\begin{itemize}
  \item des transferts de données ;
  \item des synchronisations ;
  \item des calculs de fusion ;
  \item des représentations intermédiaires ;
  \item des dépendances entre modules.
\end{itemize}

Le coût global peut donc augmenter si la communication
est trop importante.

\section{Le risque de fragmentation}

Une partition trop rigide des dimensions peut empêcher
les sondes de combiner des informations complémentaires.

Il faut préserver la possibilité de construire
des relations entre sous-espaces.

\section{Le risque de complexité du coordinateur}

Dans l'architecture hybride, le coordinateur RBF global
peut devenir le composant le plus coûteux.

Il faudra contrôler :

\[
d_g,\qquad K_g,\qquad M,\qquad d_h.
\]

Le nombre de centres et la dimension de l'espace global
devront être évalués expérimentalement.

% ============================================================
\chapter{Programme expérimental proposé}
% ============================================================

\section{Étape 1 : optimiser EX7 existant}

Avant de modifier la fonction mathématique des sondes,
il faut mesurer le coût de l'implémentation actuelle.

Travaux proposés :

\begin{itemize}
  \item profiler chaque étape d'EX7 ;
  \item mesurer le coût des encodeurs ;
  \item mesurer les distances RBF ;
  \item mesurer les exponentielles ;
  \item mesurer le coordinateur ;
  \item vérifier les possibilités de vectorisation.
\end{itemize}

\section{Étape 2 : EX7-Matrix}

Construire une variante où les sondes locales utilisent
des matrices et des activations non linéaires.

Comparer les résultats à EX7-RBF avec un budget
paramétrique contrôlé.

\section{Étape 3 : EX7-Matrix-Vector}

Conserver la même fonction mathématique qu'EX7-Matrix,
mais vectoriser les opérations.

Comparer les latences à batch égal et sur le même matériel.

\section{Étape 4 : EX7-Hybrid}

Construire la variante avec :

\[
\text{Sondes matricielles}
\longrightarrow
\text{Fusion}
\longrightarrow
\text{Coordinateur RBF}.
\]

Mesurer si la coordination RBF apporte une amélioration
réelle par rapport à une coordination matricielle.

\section{Étape 5 : spécialisation et communication}

Explorer progressivement :

\begin{itemize}
  \item les rôles distincts des sondes ;
  \item la pondération dynamique ;
  \item les sous-espaces dimensionnels ;
  \item les connexions croisées ;
  \item le passage de messages ;
  \item la coordination hiérarchique.
\end{itemize}

\section{Étape 6 : expérience linguistique}

Commencer par une tâche linguistique limitée :

\begin{itemize}
  \item classification de phrases ;
  \item identification de relations ;
  \item désambiguïsation ;
  \item prédiction de tokens ;
  \item modélisation de séquences.
\end{itemize}

Il serait prématuré de commencer directement par
un grand modèle linguistique.

\section{Tableau général des expériences}

\begin{table}[H]
\centering
\small
\begin{tabularx}{\textwidth}{
  >{\raggedright\arraybackslash}p{3.0cm}
  X
  >{\raggedright\arraybackslash}p{2.1cm}
}
\toprule
Expérience & Question & Statut\\
\midrule
EX7-RBF
&
Architecture de référence
&
Implémentée\\

EX7-Matrix
&
Les matrices accélèrent-elles les sondes ?
&
À construire\\

EX7-Matrix-Vector
&
La vectorisation réduit-elle la latence ?
&
À construire\\

EX7-Hybrid
&
La RBF globale apporte-t-elle un avantage ?
&
À construire\\

Coordination dynamique
&
Les poids doivent-ils dépendre de l'entrée ?
&
À explorer\\

Partition dimensionnelle
&
Peut-on distribuer les sous-espaces ?
&
À explorer\\

Communication entre sondes
&
Les unités peuvent-elles coopérer ?
&
À explorer\\

Modèle linguistique
&
Les sondes peuvent-elles traiter les séquences ?
&
Vision long terme\\
\bottomrule
\end{tabularx}
\caption{Programme de recherche proposé.}
\end{table}

% ============================================================
\chapter{Vision architecturale à long terme}
% ============================================================

\section{Le principe général}

La vision à long terme est celle d'une architecture
composée d'un grand nombre d'unités de traitement
relativement petites.

Chaque unité pourrait :

\begin{itemize}
  \item recevoir une partie ou une projection des données ;
  \item effectuer des transformations locales ;
  \item produire une représentation ;
  \item communiquer avec d'autres unités ;
  \item participer à une coordination globale.
\end{itemize}

La capacité du système ne serait donc pas nécessairement
concentrée dans un unique bloc dense.

Elle pourrait être distribuée entre plusieurs unités
et leurs interactions.

\section{Architecture conceptuelle générale}

\begin{figure}[H]
\centering
\begin{tikzpicture}[node distance=0.7cm]
  \node[input] (input) {
    Entrée\\
    Données ou séquence
  };

  \node[local,below=of input,text width=9cm] (local) {
    Couche de sondes\\
    Transformations locales spécialisées
  };

  \node[global,below=of local,text width=9cm] (comm) {
    Couche de communication\\
    Échanges et relations entre sondes
  };

  \node[rbf,below=of comm,text width=9cm] (coord) {
    Coordination globale\\
    Fusion, sélection et modélisation non linéaire
  };

  \node[global,below=of coord,text width=9cm] (memory) {
    Mémoire et contexte\\
    États persistants ou représentations séquentielles
  };

  \node[input,below=of memory] (output) {
    Sortie ou prédiction
  };

  \draw[flow] (input) -- (local);
  \draw[flow] (local) -- (comm);
  \draw[flow] (comm) -- (coord);
  \draw[flow] (coord) -- (memory);
  \draw[flow] (memory) -- (output);
\end{tikzpicture}
\caption{Vision générale d'une architecture distribuée à sondes.}
\end{figure}

Ce diagramme est un modèle conceptuel.
Il ne décrit pas une architecture déjà réalisée.

\section{La question de l'échelle}

Une architecture linguistique pourrait comporter :

\begin{itemize}
  \item des sondes locales ;
  \item des groupes de sondes ;
  \item des coordinateurs intermédiaires ;
  \item un ou plusieurs coordinateurs globaux ;
  \item des mécanismes de mémoire.
\end{itemize}

Une organisation hiérarchique pourrait être étudiée :

\[
\text{Sondes}
\rightarrow
\text{Groupes}
\rightarrow
\text{Coordinateurs}
\rightarrow
\text{Sortie}.
\]

Cette organisation pourrait répartir la complexité,
mais elle introduit aussi des problèmes de communication,
de stabilité et de coût.

\section{Compréhension sémantique}

Pour qu'un modèle traite la sémantique, il ne suffit pas
de répartir les mots entre plusieurs sondes.

Il doit apprendre des relations telles que :

\begin{itemize}
  \item la coréférence ;
  \item les relations sujet-verbe-objet ;
  \item les dépendances entre propositions ;
  \item les relations causales exprimées ;
  \item les ambiguïtés lexicales ;
  \item les relations entre concepts ;
  \item l'évolution du sens avec le contexte.
\end{itemize}

Une organisation spécialisée pourrait faciliter
l'analyse de ces phénomènes, à condition que les
interactions entre sondes soient suffisamment riches.

% ============================================================
\chapter{Synthèse des idées et décisions architecturales}
% ============================================================

\section{Ce qui est établi}

\begin{enumerate}
  \item EX7 V003 utilise sept sondes AdaptiveRBF.
  \item Chaque sonde reçoit les caractéristiques complètes.
  \item Les partitions concernent les exemples d'entraînement.
  \item Le coordinateur utilise des poids globaux appris.
  \item Toutes les sondes sont exécutées à l'inférence.
  \item L'architecture possède un nombre de paramètres
  comparable au MLP apparié.
  \item Sa latence mesurée est supérieure à celle du MLP.
\end{enumerate}

\section{Ce qui est proposé}

\begin{enumerate}
  \item Remplacer les sondes RBF par des sondes matricielles.
  \item Vectoriser les calculs locaux.
  \item Déplacer les fonctions RBF vers un coordinateur global.
  \item Distribuer éventuellement les dimensions entre sondes.
  \item Introduire une communication entre unités.
  \item Rendre la coordination éventuellement dynamique.
  \item Explorer des organisations hiérarchiques.
\end{enumerate}

\section{Ce qui relève de la vision à long terme}

\begin{enumerate}
  \item Construire une architecture linguistique à sondes.
  \item Étudier une spécialisation fonctionnelle des unités.
  \item Apprendre des relations sémantiques.
  \item Traiter des séquences et maintenir un contexte.
  \item Explorer des alternatives à l'attention standard.
  \item Améliorer la traçabilité des calculs internes.
  \item Étudier si la coopération de nombreuses petites unités
  peut produire des capacités complexes.
\end{enumerate}

\section{Tableau de distinction des niveaux de preuve}

\begin{table}[H]
\centering
\small
\begin{tabularx}{\textwidth}{
  >{\raggedright\arraybackslash}p{4cm}
  X
}
\toprule
Énoncé & Niveau de preuve\\
\midrule
EX7 V003 fonctionne sur le jeu de référence
&
Résultat expérimental communiqué\\

EX7 possède une latence supérieure au MLP apparié
&
Résultat expérimental communiqué\\

Les sondes matricielles seront plus rapides
&
Hypothèse à tester\\

Le coordinateur RBF global améliorera la qualité
&
Hypothèse à tester\\

Les sondes apprendront des rôles sémantiques distincts
&
Hypothèse à tester\\

L'attention pourra être remplacée
&
Hypothèse de recherche à long terme\\

Une architecture à sondes pourra traiter le langage
&
Objectif de recherche à long terme\\

La modularité facilitera l'interprétabilité
&
Hypothèse à vérifier expérimentalement\\
\bottomrule
\end{tabularx}
\caption{Séparation entre faits, hypothèses et objectifs.}
\end{table}

% ============================================================
\chapter{Conclusion générale}
% ============================================================

EX7 explore une idée architecturale : construire un système
composé d'unités de traitement relativement autonomes,
capables de produire des représentations locales et de
contribuer à une décision globale.

La première implémentation utilise des sondes RBF.
Elle démontre qu'une telle organisation peut être entraînée
sur une tâche de classification avec un budget paramétrique
proche de celui d'un petit MLP.

Les résultats disponibles sont encourageants pour poursuivre
la recherche, mais ne montrent pas une supériorité générale.
Le MLP apparié reste meilleur en log-loss et en latence
dans les mesures communiquées.

L'une des pistes les plus intéressantes est l'architecture
hybride :

\[
\boxed{
\text{Sondes matricielles locales}
+
\text{Coordination RBF globale}
}
\]

Elle permettrait de conserver le principe de modularité
tout en testant une répartition différente des fonctions
de calcul.

À plus long terme, le projet pourrait explorer une architecture
linguistique dans laquelle les sondes traitent des informations
locales, échangent des représentations, construisent des
relations et contribuent à une mémoire contextuelle.

L'attention standard n'est pas nécessairement obligatoire,
mais ses fonctions relationnelles doivent être assurées
par un mécanisme alternatif.

La question centrale du projet peut ainsi être formulée :

\begin{quote}
Peut-on construire un système d'apprentissage complexe
en répartissant les transformations, les représentations
et la coordination entre de nombreuses unités spécialisées,
tout en conservant une communication suffisante pour
apprendre des relations globales et traiter des séquences ?
\end{quote}

La réponse ne peut être obtenue par la seule conception
théorique. Elle exige une progression expérimentale,
des comparaisons contrôlées, des ablations et des validations
sur des tâches de complexité croissante.

% ============================================================
\appendix
\chapter{Sources et périmètre documentaire}
% ============================================================

\section{Sources du projet}

Les éléments décrivant l'implémentation de référence
s'appuient notamment sur les fichiers de l'archive
\texttt{sondes\_ex7\_v003.zip} :

\begin{enumerate}
  \item
  \texttt{reports/ARCHITECTURE\_V003.md} :
  description de l'organisation EX7 V003, des sondes,
  des partitions d'entraînement et de l'inférence.

  \item
  \texttt{inputs/reference\_v002/src/ex6\_common.py} :
  implémentation de la classe \texttt{AdaptiveRBF},
  encodeur, centres, largeurs et lecture radiale.
\end{enumerate}

Les chiffres expérimentaux sont repris des résultats
communiqués dans le cadre de l'analyse du projet.

Les propositions EX7-Matrix, EX7-Matrix-Vector,
EX7-Hybrid et les architectures linguistiques ne sont
pas présentées comme des implémentations déjà validées.

\section{Références conceptuelles externes}

Les mécanismes utilisés dans les propositions appartiennent
à des familles connues de l'apprentissage automatique :

\begin{itemize}
  \item réseaux à fonctions de base radiales ;
  \item perceptrons multicouches ;
  \item calcul matriciel batché ;
  \item modèles à experts et mécanismes de routage ;
  \item réseaux de neurones à passage de messages ;
  \item Transformers et attention ;
  \item modèles séquentiels et mécanismes de mémoire.
\end{itemize}

Cette note ne prétend pas établir une nouveauté
bibliographique de ces mécanismes. Une étude de l'état
de l'art sera nécessaire pour situer précisément
la contribution originale du projet.

\end{document}
